\section{Discussion}
\label{sec:discussion}

The measured matrix shows substantial task-dependent variation among the evaluated implementations. The classifier's macro-F1 is high on Banking77, CLINC150 plus, SMS Spam, and LEDGAR, and much lower on PubMedQA, VitaminC, and ESCI. Laya exceeds the classifier on PubMedQA and on Civil Comments under macro-F1, and is far from it on Banking77 and LEDGAR. The rule is close to the classifier on the only task where a rule was defined, and is not defined elsewhere. Qwen's macro-F1 is the lowest eligible score on seven of the eight tasks. On VitaminC it is 0.2209 and Laya is 0.2165. The reason a large gap appears is not the same on every task. Where \(V\) is near 0, the gap is among legal labels. Where \(V\) is near 1, the gap is the contract.

The predictor results are about a model that was not allowed to represent ``Qwen fails this kind of task'' as a separate slope. Under that restriction, the primitive mean is a strong baseline because it already carries the level of each primitive. The fingerprint sometimes moves macro-F1 error down. It usually moves invalid-rate error up, including on the tasks where the contract failure is largest. Even the helpful ESCI movement does not change the policy. LLMRouterBench reports a related pattern inside LLM routing: a method can fail to beat a simple baseline even when model behavior is heterogeneous (Li, Zhang, et al.,~\citeyear{li2026llmrouterbench}). We cite that as a caution, not as the same experiment.

We do not conclude that a fingerprint with primitive-specific slopes, a larger task sample, or a different generative configuration would fail. We also do not conclude that it would succeed. The protocol was frozen before \cref{tab:selector} existed. Changing the model to fit the miss would spend that property.
